% ─────────────────────────────────────────────────────────────────────────────
\section{Introduction and Project Context}
\label{sec:intro}
% ─────────────────────────────────────────────────────────────────────────────

This final report presents the completed architecture, implementation approach, validation
summary, and engineering assessment of the senior project titled
\textbf{Hybrid Intelligence for Deep Semantic Analysis of UAV Thermal Imagery in Closed
Networks}. The project was developed as a decision-support system for security-oriented
scenarios in which external connectivity is either prohibited or operationally unsafe.

The central engineering problem addressed by the project is the gap between
\emph{detection} and \emph{understanding}. Existing surveillance pipelines can often locate
objects such as people or vehicles, but they frequently fail to explain their likely role,
status, or threat relevance. For thermal imagery this gap is even larger because the input is
noisy, low in texture, and fundamentally different from standard RGB imagery.

The project therefore adopted a hybrid architecture: a fast detector first scans the full
thermal scene and identifies candidate objects, after which a vision-language reasoner
performs deeper analysis on cropped regions of interest. In practical terms, the system is
designed to move from outputs such as ``vehicle detected'' toward outputs such as ``light
pickup-type vehicle with an active front thermal signature'' or ``human carrying a long,
rifle-like object.''

\subsection{Project Motivation}

Thermal imagery is highly relevant for night operations, wide-area reconnaissance, and
degraded visual environments. However, most modern multimodal models are trained mainly
on RGB imagery and therefore inherit a strong bias toward color, texture, and ordinary
lighting conditions. This creates a mismatch between the strengths of general-purpose visual
models and the actual needs of thermal UAV analysis.

At the same time, security-sensitive deployments often cannot rely on cloud services. A
closed-network architecture is desirable for data sovereignty, reduced latency, and
operational robustness. These two realities together motivated a design that is both
\textbf{thermal-aware} and \textbf{offline-capable}.

\subsection{Project Objectives}

The final project was organized around the following objectives:

\begin{itemize}[nosep]
  \item detect objects of interest in thermal UAV imagery with a model specialized for fast
    scene scanning,
  \item perform deeper semantic interpretation on detected targets rather than relying only on
    coarse labels,
  \item operate fully inside a local or air-gapped network without dependence on external
    APIs or cloud inference,
  \item produce structured outputs suitable for dashboards, logs, and post-mission review,
  \item preserve a human-in-the-loop decision model rather than replacing the operator.
\end{itemize}

\subsection{Project Context}

The work aligns with the broader goal of building secure and locally deployable artificial
intelligence systems for defense and critical-infrastructure monitoring. Within this context,
the project treated the software as a decision-support mechanism, not as an autonomous
action system. The operator remains responsible for interpreting alerts and making any
mission-critical decision.

The project also fits the educational purpose of the senior design sequence: it combines
requirements engineering, system design, literature review, model selection, deployment
planning, testing, and ethical assessment within a single integrated capstone effort.

\subsection{Scope of This Report}

This report documents:

\begin{itemize}[nosep]
  \item the final problem framing and literature background,
  \item the requirements and constraints that shaped the system,
  \item the implemented two-stage architecture,
  \item the tools and technologies studied and applied,
  \item the test results and their engineering interpretation,
  \item the global, economic, environmental, societal, and ethical implications of the work.
\end{itemize}
